
\begin{table}[h!]
\centering
\resizebox{0.9\textwidth}{!}{
\begin{tabular}{l|c|c}
\toprule
\multirow{2}{*} {\bf Model} &{\bf Accuracy}& {\bf Accuracy}\\
& Pyrosetta  pre-processing& {Biopython pre-processing} \\
\midrule
\HERMESPy & 0.73 & 0.75 \\
\HERMESPyNoise & 0.64 & 0.65 \\
\HERMESPyFtRelaxed & 0.55 & 0.47 \\
\HERMESPyNoiseFtRelaxed & 0.50 & 0.44 \\
\midrule
\HERMESPyFtRos & 0.41 & 0.40 \\
\HERMESPyNoiseFtRos & 0.38 & 0.37 \\
\HERMESPyFtCdna & 0.47 & 0.45 \\
\HERMESPyNoiseFtCdna & 0.39 & 0.38 \\
\HERMESPyFtRelaxedFtCdna & 0.37 & - \\
\HERMESPyNoiseFtRelaxedFtCdna & 0.34 & - \\
\HERMESPyFtCdnaESM & 0.46 & 0.49 \\
\HERMESPyNoiseFtCdnaESM & 0.40 & 0.40 \\
\midrule
\HERMESPyUntrainedFtCdna & 0.09 & - \\
\HERMESPyNoiseUntrainedFtCdna & 0.08 & - \\
\bottomrule
\end{tabular}}
\figcaption{Accuracy of HERMES models on wildtype amino-acid classification on all sites across
40 CASP12 test proteins.}{
Accuracy is defined as proportion of sites for which the wild-type amino acid is predicted with the
highest probability among all 20 canonical amino acids.
Model names indicate the architecture, the coordinate-noise amplitude used, and when applicable, the
fine-tuning dataset (listed after ``$+$"); \textit{Untr.} is short for \textit{Untrained},
indicating models that had no pre-training and were instead only trained on stability effects.
Accuracy is reported for the two pre-processing schemes (with PyRosetta and Biopython) used in
HERMES.
}
\label{table:aa_wt_cls}
\end{table}


\begin{table}[h!]
\centering
\resizebox{0.4\textwidth}{!}{%
\begin{tabular}{l | c}
\toprule
\textbf{Model} & \textbf{hh:mm:ss} \\
\midrule
\HERMESPyNoise & 00:11:23 \\
\HERMESPyNoiseFtRelaxed & 00:11:23 \\
\HERMESPyNoiseRelaxed & 12:13:20 \\
\bottomrule
\end{tabular}}
\figcaption{Inference speed of HERMES models on the T2837 dataset.}{ Runtimes are reported in
hours (hh), minutes (mm), and seconds (ss) for inference on the T2837 dataset, which comprises 2837
mutation effects across 129 proteins. For HERMES-{\em fixed} and HERMES-{\em amortized}, the script
`mutation\_effect\_prediction\_with\_hermes.py` was used; for HERMES-{\em relaxed}, the script
`mutation\_effect\_prediction\_with\_hermes\_with\_relaxation.py` was used. Both scripts, along with
the dataset `csv` file, are available in our GitHub repository. All models were executed using a
single CPU and a single A40 GPU.}
\label{table:inference_speed}
\end{table}


\begin{table}[h!]
\centering
\footnotesize
\setlength{\tabcolsep}{6pt}
\begin{tabularx}{\textwidth}{>{\raggedright\arraybackslash}p{0.27\linewidth}
>{\raggedright\arraybackslash}X}
\toprule
\textbf{Category} & \textbf{Description} \\
\midrule

\multirow{18}{*}{Hydrophobic Property}
& 1.\ Retention coefficient in TFA \\
& 2.\ Free energy of solution in water \\
& 3.\ Solvation free energy \\
& 4.\ Melting point \\
& 5.\ Number of hydrogen-bond donors \\
& 6.\ Number of full nonbonding orbitals \\
& 7.\ Partition energy \\
& 8.\ Hydration number \\
& 9.\ Retention coefficient in high performance liquid chromatography (HPLC), pH 7.4 \\
& 10.\ Retention coefficient in HPLC, pH 2.1 \\
& 11.\ Partition coefficient in thin-layer chromatography \\
& 12.\ Retention coefficient at pH 2 \\
& 13.\ $R_f$ for 1-N-(4-nitrobenzofurazono)-amino acids in ethyl acetate/pyridine/water \\
& 14.\ $\Delta G$ of transfer from organic solvent to water \\
& 15.\ Hydration potential or free energy of transfer from vapor phase to water \\
& 16.\ $R_f$, salt chromatography \\
& 17.\ $\log D$, partition coefficient at pH 7.1 for acetamide derivatives of amino acids in
octanol/water \\
& 18.\ $\Delta G = RT \log f$, $f$ = fraction buried/accessible amino acids in 22 proteins \\
\midrule

\multirow{17}{*}{Steric Property}
& 19.\ Average volume of buried residue \\
& 20.\ Residue accessible surface area in tripeptide \\
& 21.\ Graph shape index \\
& 22.\ Normalized van der Waals volume \\
& 23.\ STERMIMOL length of the side chain \\
& 24.\ STERMIMOL minimum width of the side chain \\
& 25.\ STERMIMOL maximum width of the side chain \\
& 26.\ Average accessible surface area \\
& 27.\ Distance between $C_{\alpha}$ and centroid of side chain \\
& 28.\ Side-chain angle $\theta$ \\
& 29.\ Side-chain torsion angle $\phi$ \\
& 30.\ Radius of gyration of side chain \\
& 31.\ Van der Waals parameter $R_0$ \\
& 32.\ Van der Waals parameter $\varepsilon$ \\
& 33.\ Refractivity \\
& 34.\ Value of $\theta$ (i) \\
& 35.\ Substituent van der Waals volume \\
\midrule

\multirow{15}{*}{Electronic Property}
& 36.\ $\alpha$CH chemical shifts \\
& 37.\ $\alpha$NH chemical shifts \\
& 38.\ A parameter of charge transfer capability \\
& 39.\ A parameter of charge transfer donor capability \\
& 40.\ Nuclear magnetic resonance (NMR) chemical shift of $\alpha$ carbon \\
& 41.\ Localized electrical effect \\
& 42.\ Positive charge \\
& 43.\ Negative charge \\
& 44.\ Polarity \\
& 45.\ Net charge \\
& 46.\ Amphipathicity index \\
& 47.\ Isoelectric point \\
& 48.\ Electron-ion interaction potential values \\
& 49.\ $\mathrm{pK}_{\mathrm{NH}_2}$ (NH$_2$ on $C_{\alpha}$) \\
& 50.\ $\mathrm{pK}_{\mathrm{COOH}}$ (COOH on $C_{\alpha}$) \\
\bottomrule
\end{tabularx}
\figcaption{Table of amino-acid properties used for comparison with substitution matrices.}{
Values associated with each amino acid are listed in ref.~\cite{mei_new_2005}.}
\label{table:aa_properties}
\end{table}


\begin{table}[h!]
    \centering
    \resizebox{\textwidth}{!}{
        \begin{tabular}{c c c c | c c c c c c c}
            \toprule
            \multirow{2}{*}{\makecell{antigen\\PDB id}} &
            \multirow{2}{*}{\makecell{stabilized\\name}} &
            \multirow{2}{*}{mutation} &
            \multirow{2}{*}{\makecell{BLOSUM62\\score ; rank}} &
                \makecell{Rosetta} & \makecell{Protein\\MPNN\\0.30} & \makecell{Thermo\\MPNN\\+
                Megascale} & \makecell{HERMES\\-{\em fixed} 0.50} & \makecell{HERMES\\-{\em
                amortized}\\0.50} & \makecell{HERMES\\-{\em fixed} 0.50\\+ cDNA117k} &
                \makecell{HERMES\\-{\em fixed} 0.50\\+ Megascale} \\
                \cmidrule(lr){5-11}
                & & & & $r_\wt \rightarrow r_\mt$ & $r_\wt \rightarrow r_\mt$ & $r_\wt \rightarrow
                r_\mt$ & $r_\wt \rightarrow r_\mt$ & $r_\wt \rightarrow r_\mt$ & $r_\wt \rightarrow
                r_\mt$ & $r_\wt \rightarrow r_\mt$ \\
                \midrule
            \multirow{7}{*}{\makecell{RSV-F\\4JHW}} & \multirow{2}{*}{Cav1
            \citep{mclellan_structure-based_2013}} & S190F & -2 ; 16 & \yellowl 19 $\rightarrow$ 11
            & 4 $\rightarrow$ 7 & \green 14 $\rightarrow$ 3 & 4 $\rightarrow$ 8 & \greenl 11
            $\rightarrow$ 4 & \green 12 $\rightarrow$ 3 & \green 13 $\rightarrow$ 3 \\
             & & V207L & 1 ; 3 & \greenl 9 $\rightarrow$ 4 & 2 $\rightarrow$ 7 & \greenl 6
             $\rightarrow$ 5 & 2 $\rightarrow$ 3 & \green 3 $\rightarrow$ 2 & \green 3 $\rightarrow$
             2 & \green 3 $\rightarrow$ 2 \\
            \cmidrule(lr){2-11}
             & \multirow{1}{*}{Uncl. \citep{lee_rational_2024}} & S215P & -1 ; 13 & \green 13
             $\rightarrow$ 1 & 1 $\rightarrow$ 14 & 10 $\rightarrow$ 20 & 10 $\rightarrow$ 15 &
             \green 6 $\rightarrow$ 3 & 16 $\rightarrow$ 19 & 4 $\rightarrow$ 20 \\
            \cmidrule(lr){2-11}
             & \multirow{4}{*}{TriC \citep{mclellan_structure-based_2013}} & D486H & -1 ; 7 & \green
             9 $\rightarrow$ 2 & 1 $\rightarrow$ 10 & 1 $\rightarrow$ 8 & \greenl 5 $\rightarrow$ 4
             & \greenl 10 $\rightarrow$ 4 & \yellowl 20 $\rightarrow$ 9 & \yellowl 18 $\rightarrow$
             9 \\
             & & E487Q & 2 ; 3 & 2 $\rightarrow$ 9 & 1 $\rightarrow$ 13 & 1 $\rightarrow$ 12 & 2
             $\rightarrow$ 3 & 4 $\rightarrow$ 5 & 11 $\rightarrow$ 12 & \yellowl 14 $\rightarrow$
             12 \\
             & & F488W & 1 ; 3 & 15 $\rightarrow$ 17 & 1 $\rightarrow$ 14 & \green 3 $\rightarrow$ 1
             & 2 $\rightarrow$ 6 & 1 $\rightarrow$ 13 & 1 $\rightarrow$ 3 & 1 $\rightarrow$ 2 \\
             & & D489H & -1 ; 7 & \green 5 $\rightarrow$ 2 & 3 $\rightarrow$ 14 & 5 $\rightarrow$ 12
             & \yellowl 11 $\rightarrow$ 8 & \yellowl 18 $\rightarrow$ 15 & \yellowl 20
             $\rightarrow$ 11 & \yellowl 17 $\rightarrow$ 12 \\
            \midrule
            \multirow{3}{*}{\makecell{HA\\7VDF}} & \multirow{3}{*}{Universal-HA
            \citep{milder_universal_2022}} & H355W & -2 ; 12 & \green 5 $\rightarrow$ 1 & \green 5
            $\rightarrow$ 1 & \green 4 $\rightarrow$ 2 & \green 4 $\rightarrow$ 3 & \green 7
            $\rightarrow$ 2 & \green 5 $\rightarrow$ 3 & \green 8 $\rightarrow$ 3 \\
             & & K380I & -3 ; 19 & \green 17 $\rightarrow$ 1 & \green 9 $\rightarrow$ 1 & \green 11
             $\rightarrow$ 1 & \greenl 8 $\rightarrow$ 4 & \green 14 $\rightarrow$ 3 & \green 13
             $\rightarrow$ 2 & \green 17 $\rightarrow$ 2 \\
             & & E432I & -3 ; 19 & \yellowl 14 $\rightarrow$ 9 & \green 12 $\rightarrow$ 2 & \green
             13 $\rightarrow$ 2 & \greenl 9 $\rightarrow$ 4 & \greenl 15 $\rightarrow$ 5 & \green 12
             $\rightarrow$ 3 & \green 14 $\rightarrow$ 3 \\
            \midrule
            \multirow{11}{*}{\makecell{hMPV-F\\5WB0}} & \multirow{5}{*}{M104
            \citep{gonzalez_general_2024}} & L130D & -4 ; 19 & \green 15 $\rightarrow$ 3$^\dagger$ &
            8 $\rightarrow$ 9 & 4 $\rightarrow$ 10 & 5 $\rightarrow$ 9 & \green 12 $\rightarrow$ 2 &
            \green 8 $\rightarrow$ 2 & \greenl 10 $\rightarrow$ 4 \\
             & & A159L & -1 ; 10 & \green 6 $\rightarrow$ 1$^\dagger$ & 1 $\rightarrow$ 12 & 3
             $\rightarrow$ 5 & 1 $\rightarrow$ 7 & 2 $\rightarrow$ 4 & \green 7 $\rightarrow$ 3 &
             \green 5 $\rightarrow$ 3 \\
             & & V203I & 3 ; 2 & \green 4 $\rightarrow$ 3$^\dagger$ & 1 $\rightarrow$ 2 & 1
             $\rightarrow$ 2 & 1 $\rightarrow$ 2 & \green 2 $\rightarrow$ 1 & \green 3 $\rightarrow$
             1 & \green 2 $\rightarrow$ 1 \\
             & & V430Q & -2 ; 10 & \green 12 $\rightarrow$ 3$^\dagger$ & \greenl 5 $\rightarrow$ 4 &
             \greenl 12 $\rightarrow$ 6 & \green 8 $\rightarrow$ 1 & \greenl 10 $\rightarrow$ 5 & 5
             $\rightarrow$ 12 & \yellowl 10 $\rightarrow$ 8 \\
             & & V449D & -3 ; 16 & \yellowl 13 $\rightarrow$ 8$^\dagger$ & \green 13 $\rightarrow$ 3
             & \green 16 $\rightarrow$ 2 & \green 11 $\rightarrow$ 2 & \green 15 $\rightarrow$ 2 &
             \green 18 $\rightarrow$ 2 & \greenl 19 $\rightarrow$ 6 \\
            \cmidrule(lr){2-11}
             & \multirow{4}{*}{MPV-2cREKR \citep{bakkers_efficacious_2024}} & V112R & -3 ; 19 & 6
             $\rightarrow$ 13$^\dagger$ & \green 8 $\rightarrow$ 1 & \green 2 $\rightarrow$ 1 & 2
             $\rightarrow$ 9 & \green 8 $\rightarrow$ 3 & 4 $\rightarrow$ 11 & 5 $\rightarrow$ 10 \\
             & & D209E & 2 ; 2 & \green 8 $\rightarrow$ 1$^\dagger$ & \green 11 $\rightarrow$ 1 &
             \greenl 12 $\rightarrow$ 4 & \green 3 $\rightarrow$ 1 & \green 14 $\rightarrow$ 3 &
             \greenl 15 $\rightarrow$ 6 & \yellowl 16 $\rightarrow$ 7 \\
             & & V231I & 3 ; 2 & \green 3 $\rightarrow$ 1$^\dagger$ & \green 3 $\rightarrow$ 1 &
             \green 3 $\rightarrow$ 1 & \green 2 $\rightarrow$ 1 & \green 2 $\rightarrow$ 1 & \green
             3 $\rightarrow$ 1 & \green 3 $\rightarrow$ 1 \\
             & & E453P & -1 ; 10 & \green 10 $\rightarrow$ 2$^\dagger$ & 1 $\rightarrow$ 20 & 1
             $\rightarrow$ 2 & 8 $\rightarrow$ 19 & \green 13 $\rightarrow$ 1 & \yellowl 17
             $\rightarrow$ 16 & \greenl 19 $\rightarrow$ 6 \\
            \cmidrule(lr){2-11}
             & \multirow{2}{*}{Uncl. \citep{lee_rational_2024}} & E80D & 2 ; 2 & \yellowl 18
             $\rightarrow$ 13 & \green 3 $\rightarrow$ 1 & 13 $\rightarrow$ 18 & 1 $\rightarrow$ 12
             & 7 $\rightarrow$ 17 & 15 $\rightarrow$ 19 & 14 $\rightarrow$ 18 \\
             & & V155P & -2 ; 13 & 10 $\rightarrow$ 20 & 4 $\rightarrow$ 20 & 1 $\rightarrow$ 20 & 1
             $\rightarrow$ 19 & 1 $\rightarrow$ 20 & 2 $\rightarrow$ 20 & 4 $\rightarrow$ 20 \\
            \midrule
            \multirow{8}{*}{\makecell{DENV-E\\1OAN}} & \multirow{8}{*}{SC12
            \citep{phan_conserved_2022}} & S29K & 0 ; 6 & \green 4 $\rightarrow$ 3$^\dagger$ &
            \green 3 $\rightarrow$ 1 & \green 7 $\rightarrow$ 2 & 1 $\rightarrow$ 12 & 2
            $\rightarrow$ 5 & 6 $\rightarrow$ 12 & 9 $\rightarrow$ 14 \\
             & & T33V & 0 ; 3 & \greenl 9 $\rightarrow$ 4$^\dagger$ & \green 3 $\rightarrow$ 1 &
             \green 6 $\rightarrow$ 1 & \green 3 $\rightarrow$ 1 & \green 5 $\rightarrow$ 1 & \green
             10 $\rightarrow$ 1 & \green 10 $\rightarrow$ 1 \\
             & & A35M & -1 ; 13 & \greenl 11 $\rightarrow$ 4$^\dagger$ & \green 6 $\rightarrow$ 1 &
             \yellowl 10 $\rightarrow$ 8 & 1 $\rightarrow$ 12 & 1 $\rightarrow$ 11 & 1 $\rightarrow$
             7 & 6 $\rightarrow$ 7 \\
             & & G106D & -1 ; 5 & \green 18 $\rightarrow$ 1$^\dagger$ & \green 11 $\rightarrow$ 1 &
             6 $\rightarrow$ 11 & \green 8 $\rightarrow$ 1 & \green 6 $\rightarrow$ 2 & \green 20
             $\rightarrow$ 1 & \green 18 $\rightarrow$ 1 \\
             & & A259W & -3 ; 20 & \green 7 $\rightarrow$ 1$^\dagger$ & 1 $\rightarrow$ 3 & 4
             $\rightarrow$ 11 & 1 $\rightarrow$ 20 & 1 $\rightarrow$ 13 & 1 $\rightarrow$ 2 & \green
             2 $\rightarrow$ 1 \\
             & & T262R & -1 ; 11 & \green 15 $\rightarrow$ 3$^\dagger$ & 1 $\rightarrow$ 15 &
             \yellowl 18 $\rightarrow$ 7 & 3 $\rightarrow$ 14 & \green 10 $\rightarrow$ 3 & \green
             16 $\rightarrow$ 3 & \greenl 17 $\rightarrow$ 4 \\
             & & F279W & 1 ; 3 & \green 3 $\rightarrow$ 1$^\dagger$ & 4 $\rightarrow$ 11 & \green 4
             $\rightarrow$ 1 & 1 $\rightarrow$ 6 & 3 $\rightarrow$ 7 & 1 $\rightarrow$ 4 & 1
             $\rightarrow$ 2 \\
             & & T280P & -1 ; 13 & \green 12 $\rightarrow$ 2$^\dagger$ & 7 $\rightarrow$ 10 & 9
             $\rightarrow$ 14 & 4 $\rightarrow$ 19 & \green 6 $\rightarrow$ 2 & 2 $\rightarrow$ 20 &
             8 $\rightarrow$ 19 \\
            \midrule
            \multirow{4}{*}{\makecell{SARS-Cov-2\\6VSB}} & \multirow{4}{*}{hexapro
            \citep{hsieh_structure-based_2020}} & F817P & -4 ; 20 & \green 3 $\rightarrow$ 1 &
            \green 16 $\rightarrow$ 1 & 1 $\rightarrow$ 12 & 2 $\rightarrow$ 16 & \green 3
            $\rightarrow$ 1 & 2 $\rightarrow$ 18 & 2 $\rightarrow$ 19 \\
             & & A892P & -1 ; 14 & \green 16 $\rightarrow$ 3 & \green 5 $\rightarrow$ 1 & 2
             $\rightarrow$ 4 & \green 2 $\rightarrow$ 1 & \green 4 $\rightarrow$ 1 & \green 4
             $\rightarrow$ 1 & \green 5 $\rightarrow$ 1 \\
             & & A899P & -1 ; 14 & 8 $\rightarrow$ 19 & 6 $\rightarrow$ 8 & 7 $\rightarrow$ 18 & 2
             $\rightarrow$ 6 & \green 4 $\rightarrow$ 1 & 8 $\rightarrow$ 15 & 5 $\rightarrow$ 17 \\
             & & A942P & -1 ; 14 & 5 $\rightarrow$ 12 & \green 3 $\rightarrow$ 1 & 2 $\rightarrow$ 8
             & 1 $\rightarrow$ 4 & \green 2 $\rightarrow$ 1 & \green 2 $\rightarrow$ 1 & \green 2
             $\rightarrow$ 1 \\
            \midrule
            & \multicolumn{3}{c|}{\green \makecell{proportion of correctly and\\strongly suggested
            mutations}} & \textbf{20/33} & \textbf{15/33} & \textbf{11/33} & 8/33 & \textbf{19/33} &
            \textbf{15/33} & \textbf{13/33} \\
            \cmidrule(lr){2-11}
            & \multicolumn{3}{c|}{\greenl \makecell{proportion of correctly and\\at least moderately
            suggested mutations}} & \textbf{23/33} & \textbf{16/33} & 14/33 & 11/33 & \textbf{23/33}
            & \textbf{16/33} & \textbf{17/33} \\
            \cmidrule(lr){2-11}
            & \multicolumn{3}{c|}{\makecell{proportion of correctly and\\at least weakly suggested
            mutations}} & \textbf{27/33} & 16/33 & 16/33 & 12/33 & \textbf{24/33} & 19/33 & 22/33 \\

            \bottomrule
        \end{tabular}
    }
    \figcaption{Predicting antigen-stabilizing mutations with HERMES: extended results.}{ Recall for different
models (columns) is evaluated on 33 previously reported antigen-stabilizing mutations (rows)
spanning five viral antigens. For each antigen, we list the PDB structure used for scoring and the
publication(s) that originally reported the mutation. Mutations are specified as
wild-type$\to$mutant substitutions at the annotated site. Seven models are compared (columns). We
additionally report the BLOSUM62 substitution score for each mutation and the mutant's rank among
the 20 possible amino-acid substitutions for the wild-type residue (per BLOSUM62). For each model
and mutation, predicted ranks of the wild-type and mutant amino acids are shown as $r_\wt \to
r_\mt$. Dagger symbols ($\dagger$) indicate mutations originally proposed as stabilizing by
Rosetta-based pipelines in the source reference. When a model ranks the mutant better than the wild
type ($r_\mt < r_\wt$), the cell is shaded by the prediction strength based on the value of $r_\mt$:
dark green, strongly suggested ($r_\mt \leq 3$); light green, moderately suggested ($4 \leq r_\mt
\leq 6$); light yellow, weakly suggested ($r_\mt \geq 6$). Column summaries report counts of
strongly, at least moderately, and at least weakly suggested mutations (out of 33); {\bf bold}
indicates significance (p-value $<0.05$) for the number of recalled mutations relative to a random
null model (see Fig.~\ref{fig:antigen_pvalues_vs_random} for p-values and Methods for details.) All
structures were scored in their native multimeric states, generating symmetric partners when needed.
ThermoMPNN's native mode predicts mutation effects only for monomers, ignoring multimeric assemblies
even when present in the input structure. ``Uncl." stands for uncleaved prefusion-closed state.
    }
    \label{table:antigen_results}
\end{table}



\begin{table}[h!]
    \centering
    \begin{tabular}{c c | c c | c c}
        \toprule
          & & \makecell{\textbf{HERMES GPU}\\{[seconds]}} & \makecell{\textbf{HERMES
          CPU}\\{[seconds]}} & \multicolumn{2}{|c}{\makecell{\textbf{Rosetta}\\{[CPU-hours]}}} \\
        pdbID & \makecell{\# of monomer sites} & \makecell{all sites} & \makecell{all sites} &
        \makecell{one site} & \makecell{all sites} \\
        \midrule
        4JHW & 449 & 57 & 112 & 150 & 67,350 \\
        7VDF & 485 & 64 & 118 & 164 & 79,540 \\
        5WB0 & 442 & 43 & 154 & 151 & 66,742 \\
        1OAN & 394 & 32 & 151 & 59 & 23,246 \\
        6VSB & 968 & 69 & 265 & 156 & 151,008 \\
        \bottomrule
    \end{tabular}
    \figcaption{Execution times for saturation mutagenesis predictions on the viral antigens considered
    in this study.}{ Executions times (in seconds) of HERMES apply to HERMES-{\em fixed} and
    HERMES-{\em amortized} models, regardless of whether zero-shot or fine-tuned. Times were
    computed when running the script \texttt{run\_hermes\_on\_pdbfiles.py} providing as input the
    pdbfile as well as a single monomeric chain. A single CPU core with 64GB of memory was used, and
    a NVIDIA A40 GPU when applicable. For Rosetta, we computed times (in CPU-hours) for a single CPU
    core with 4 GBs of memory, and averaging 10 relaxation instances, which we consider the minimum
    number of instances for robust results. Times for all sites in the structure were extrapolated
    by multiplying the calculated average time for a single mutation by the number of monomeric
    sites.
    }
    \label{table:speed_on_antigens}
\end{table}

\begin{table}[h!]
    \centering
    \resizebox{\textwidth}{!}{
        \begin{tabular}{c c c | c c c}
            \toprule
            \makecell{antigen\\PDB id} &
            \makecell{stabilized\\name} &
            \makecell{mutation} &
            \makecell{mutation\\type} &
            \makecell{is synergistic} &
            \makecell{notes} \\

                \midrule
            \multirow{7}{*}{\makecell{RSV-F\\4JHW}} & \multirow{2}{*}{Cav1
            \citep{mclellan_structure-based_2013}} & S190F & cavity-filling & False & \\
             &  & V207L & cavity-filling & False &   \\
            \cmidrule(lr){2-6}
             & \multirow{1}{*}{Uncl. \citep{lee_rational_2024}} & S215P & proline & False &   \\
            \cmidrule(lr){2-6}
             & \multirow{4}{*}{TriC \citep{mclellan_structure-based_2013}} & D486H & electrostatic &
             True & \\
             &  & E487Q & electrostatic & True &   \\
             &  & F488W & cavity-filling & True &   \\
             &  & D489H & electrostatic & True &   \\
            \midrule
            \multirow{3}{*}{\makecell{HA\\7VDF}} & \multirow{3}{*}{Universal-HA
            \citep{milder_universal_2022}} & H355W & cavity-filling & False & \\
             &  & K380I & cavity-filling & False &   \\
             &  & E432I & cavity-filling & False &   \\
            \midrule
            \multirow{11}{*}{\makecell{hMPV-F\\5WB0}} & \multirow{5}{*}{M104
            \citep{gonzalez_general_2024}} & L130D & electrostatic & False & \\
             &  & A159L & cavity-filling & False &   \\
             &  & V203I & cavity-filling & False &   \\
             &  & V430Q & electrostatic & False &   \\
             &  & V449D & electrostatic & False &   \\
            \cmidrule(lr){2-6}
             & \multirow{4}{*}{MPV-2cREKR \citep{bakkers_efficacious_2024}} & V112R & electrostatic
             & False & \\
             &  & D209E &  & False & same charge, slightly different size: unclear  \\
             &  & V231I & cavity-filling & False &   \\
             &  & E453P & proline & False &   \\
            \cmidrule(lr){2-6}
             & \multirow{2}{*}{Uncl. \citep{lee_rational_2024}} & E80D & & False & same charge,
             slightly different size: unclear \\
             &  & V155P & proline & False &   \\
            \midrule
            \multirow{8}{*}{\makecell{DENV-E\\1OAN}} & \multirow{8}{*}{SC12
            \citep{phan_conserved_2022}} & S29K & electrostatic & False & \\
             &  & T33V & cavity-filling & False &   \\
             &  & A35M & cavity-filling & False &   \\
             &  & G106D & electrostatic & False &   \\
             &  & A259W & cavity-filling & True &   \\
             &  & T262R & electrostatic & True &   \\
             &  & F279W & cavity-filling & False &   \\
             &  & T280P & proline & False &   \\
            \midrule
            \multirow{4}{*}{\makecell{SARS-Cov-2\\6VSB}} & \multirow{4}{*}{hexapro
            \citep{hsieh_structure-based_2020}} & F817P & proline & False & \\
             &  & A892P & proline & False &   \\
             &  & A899P & proline & False &   \\
             &  & A942P & proline & False &   \\

            \bottomrule
        \end{tabular}
    }
    \figcaption{Characteristics of antigen-stabilizing mutations.}{
    Hand-curated mutation types are listed for antigen-stabilizing mutations reported in
    Fig.~\ref{fig:antigen_figure_main_results} and Table~\ref{table:antigen_results}. Cavity-filling
    mutations are defined as substitutions to hydrophobic residues that are larger than the
    wild-type when the wild-type is also hydrophobic. Electrostatic mutations are substitutions that
    change the residue's net charge. Proline mutations correspond to substitutions to proline.
    Synergistic mutations were identified through structural reasoning based on the spatial
    arrangement of mutations within the corresponding structure; see
    ref.~\citep{byrne_principles_2022} for a breakdown of mutation types considered in the
    structure-based vaccine design literature. ``Uncl." stands for ``Uncleaved Prefusion-Closed".
     }
    \label{table:mutation_types}
\end{table}


\clearpage

\begin{figure}[h!]
    \centering
    \includegraphics[width=1.0\textwidth]{hermes/figures/zeroshot_models_against_each_other.png}
    \figcaption{Comparison of zero-shot model predictions on the Megascale test set.}{ For each
    substitution in the Megascale test set, the predicted change in amino acid propensity upon
    mutation ($\delta \log p$) is compared between two models in each panel. Color indicates local
    point density (blue denotes low density and yellow denotes high density). The reported “Pr” in
    each panel corresponds to the Pearson correlation coefficient between the predictions of the
    model pair. Model names indicate the architecture and the coordinate-noise amplitude used.
    $\Delta \hat{F}$ indicates the model's prediction, following
    Equations~\ref{eq:log_ratio_wt_and_mt}~and~\ref{eq:deltaF_variants}.}
    \label{fig:relaxed_vs_ft_relaxed}
\end{figure}


\begin{figure}[ht!]
    \centering
    \includegraphics[width=0.85\textwidth]{hermes/figures/radial_plots__no_numbers.pdf}
    \figcaption{Predicting mutational effects on thermodynamic folding stability.}{
    Stabilizing-versus-destabilizing classification metrics are computed using $\Delta\Delta G < 0$
    (experimental) and $\Delta \log p > 0$ (predicted) as cutoffs for stabilizing mutations.
    \textbf{(A)} Evaluation on the T2837 results: zero-shot models (top) and models fine-tuned on
    cDNA117k (bottom). \textbf{(B)} Evaluation on Megascale test set results: zero-shot models (top)
    and models fine-tuned on the Megascale training set (bottom). Model names indicate the
    architecture, the coordinate-noise amplitude used, and when applicable, the fine-tuning dataset
    (listed after ``$+$"); \textit{Untr.} is short for \textit{Untrained}, indicating models that
    had no pre-training and were instead only trained on stability effects. Only models trained
    without coordinate noise are shown; the noise amplitude is indicated within each model name.
    }
   \label{fig:radial_plots__no_noise} 
\end{figure}

\begin{figure}
    \centering
    \includegraphics[width=1.0\linewidth]{hermes/figures/pairwise_pvalues_permutation_bucketed_by_sizereduced_precision_recall_f1.pdf}
    \figcaption{Statistical significance of model performance differences on stabilizing mutation
    identification.}{ Shown are two-tailed p-values for differences in model performance on
    identifying stabilizing mutations across Megascale test set subsets. P-values correspond to the
    performance comparisons shown in
    Fig.~\ref{fig:megascale_precision_recall_f1_by_size_cutoff_0p0}. Green indicates strong
    statistical significance ($p < 0.01$), while yellow indicates weaker significance ($0.01 \leq p
    < 0.05$). P-values were computed using a permutation test and corrected for multiple comparisons
    using the Holm–Bonferroni procedure within each performance metric (see Methods for details).}
\label{fig:pairwise_pvalues_permutation_bucketed_by_sizereduced_precision_recall_f1}
\end{figure}

\begin{figure}
    \centering
    \includegraphics[width=0.4\linewidth]{hermes/figures/pvalues__bucketed_by_size__between_small_large_buckets__vertical.pdf}
    \figcaption{Statistical significance of within-model performance differences in identifying
    stabilizing mutations across mutational size classes.}{ Shown are two-tailed P-values for
    within-model differences in performance when identifying stabilizing mutations from the
    small$\to$large vs. large$\to$small mutational subsets of the Megascale test set. P-values
    correspond to the performance comparisons shown in
    Fig.~\ref{fig:megascale_precision_recall_f1_by_size_cutoff_0p0}. Green indicates strong
    statistical significance ($p < 0.01$), while yellow indicates weaker significance ($0.01\leq p <
    0.05$). P-values were computed using a bootstrap test and corrected for multiple comparisons
    using the Holm–Bonferroni procedure within each performance metric (see Methods for details).}
    \label{fig:pvalues__bucketed_by_size__between_small_large_buckets__vertical}
\end{figure}

\begin{figure}
    \centering
    \includegraphics[width=1.0\linewidth]{hermes/figures/t2837_and_megascale_pairwise_pvalues_permutation.pdf}
    \figcaption{Statistical significance of model performance differences in identifying stabilizing
    mutations on the full test set.}{
     Shown are two-tailed p-values for differences in model performance when identifying stabilizing
     mutations on the full test set. P-values correspond to the performance comparisons shown in
     Fig.~\ref{fig:radial_plots__noise}. Green indicates strong statistical significance ($p <
     0.01$), while yellow indicates weaker significance ($0.01\leq p < 0.05$). P-values were
     computed using a permutation test and corrected for multiple comparisons using the
     Holm–Bonferroni procedure within each performance metric (see Methods for details).}
    \label{fig:t2837_and_megascale_pairwise_pvalues_permutation}
\end{figure}

\begin{figure}
    \centering
    \includegraphics[width=0.7\linewidth]{hermes/figures/ddg_stability_single_structure_results_with_bp.pdf}
    \figcaption{Pearson correlation between model predictions and experimental stability
    effects on the RaSP test set (8 proteins)~\citep{blaabjerg_rapid_2023}.}{
    Each dot represents one protein, and the horizontal bar indicates the mean correlation across
    proteins. Model labels specify the architecture, the coordinate-noise amplitude, and, when
    applicable, the fine-tuning dataset (denoted after ``$+$"). ``Bp" indicates the use of our
    open-source Biopython-based protein Pre-processing. See Methods for details on the RaSP
    dataset.}
    \label{fig:rasp_exp}
\end{figure}

\begin{figure}
    \centering
    \includegraphics[width=0.9\linewidth]{hermes/figures/t2837_pearsonr_broken_down.pdf}
    \figcaption{Pearson correlation between model predictions and experimental stability effects
    on the T2837 dataset and its subsets.}{
Pearson correlation values for all models other than HERMES are taken
from~\citep{diaz_stability_2024}. This figure closely replicates a figure from
ref.~\citep{diaz_stability_2024}, with the key difference that predictions for ``reverse" mutations
are computed here by conditioning on wild-type structures (denoted as ``wt rev"). This distinction
is made to avoid confusion with ``reverse" mutation predictions computed on mutant structures in the
Ssym dataset (Fig.~\ref{fig:ssym_figure}). For each dataset (x-axis), we denote in text the
performance of the HERMES models as well as that of the best-performing model overall.}
    \label{fig:t2837_broken_down_and_comparison}
\end{figure}

\begin{figure}
    \centering
    \includegraphics[width=0.5\linewidth]{hermes/figures/radial_plots__esmfold.pdf}
    \figcaption{Predicting mutational effects on thermodynamic stability using ESMfold
    predicted structures for fine-tuning or testing.}{ Stabilizing-versus-destabilizing
    classification metrics are computed using $\Delta\Delta G < 0$ (experimental) and $\Delta \log p
    > 0$ (predicted) as cutoffs for stabilizing mutations. We report results on the T2837 dataset,
    after fine-tuning models on cDNA117k. We consider models fine-tuned and evaluated on crystal
    structures (purple), models fine-tuned on ESMfold predicted structures and evaluated on crystal
    structures (grey, ``\textit{train} ESMfold" in the model name), and models fine-tuned on crystal
    structures and evaluated on ESMfold predicted structures (olive, ``\textit{eval} ESMfold" in the
    model name).
    }
    \label{fig:radial_plots__esmfold}
\end{figure}



\begin{figure}
    \centering
    \includegraphics[width=0.95\textwidth]{hermes/figures/mean_of_abs_symmetric__significance_matrices_all_models_vs_aminoacid_properties__holm_within.pdf}
    \figcaption{Significance testing results for
    Fig.~\ref{fig:correlations_to_aa_properties}B.}{
    Shown are p-values of two-tailed t-tests comparing the distributions of spearman correlations
    between each model's substitution matrix, and amino-acid properties of a particular class
    (electronic, hydrophobic, steric). The Holm-Bonferroni method was used to correct p-values for
    multiple testing error. Entries corresponding to pairs of distributions with p-value $\leq 0.01$
    are colored in green, p-value $\leq 0.05$ are colored in yellow, and p-value $> 0.05$ are
    colored in gray.}
    \label{fig:correlation_to_aa_properties__significance}
\end{figure}

\begin{figure}
    \centering
    \includegraphics[width=0.6\linewidth]{hermes/figures/mean_of_abs_symmetric_all__average_prediction_matrices__spearman_matrix.pdf}
    \figcaption{Spearman correlations between pairs of model-average substitution matrices
    $M^{\model}$.}{ The heatmaps for the underlying model-predicted substitution matrices are shown
    in Fig.~\ref{fig:correlations_to_aa_properties}.}
    \label{fig:spearmanr_between_model_average_substitution_matrices}
\end{figure}

\begin{figure}
    \centering
    \includegraphics[width=0.8\textwidth]{hermes/figures/mean_of_abs_symmetric__matrix_heatmaps__by_loc__part_1.pdf}
    \figcaption{Model-averaged substitution matrices stratified by protein core and surface
    residues for zero-shot models.}{
Shown are model-averaged substitution matrices $M^{\model}$ for different zero-shot models (rows
2-4), computed from subsets of sites in the Megascale test set. The first row shows the experimental
matrices for mean $|\Delta \Delta G|$ values across mutation subsets. Columns correspond to all
residues (left), core residues with solvent-accessible surface area SASA $< 1 \AA^2$ (center), and
surface residues with SASA $> 3 \AA^2$ (right).}
    \label{fig:average_prediction_matrices__abs_symm__part_1}
\end{figure}

\begin{figure}
    \centering
    \includegraphics[width=0.8\textwidth]{hermes/figures/mean_of_abs_symmetric__matrix_heatmaps__by_loc__part_2.pdf}
    \figcaption{Model-averaged substitution matrices stratified by protein core and surface
    residues for stability fine-tuned models.}{
     Similar to Fig.~\ref{fig:average_prediction_matrices__abs_symm__part_1} but for models
     fine-tuned on the Megascale dataset.}
    \label{fig:average_prediction_matrices__abs_symm__part_2}
\end{figure}

\begin{figure}
    \centering
    \includegraphics[width=0.9\linewidth]{hermes/figures/ssym_antisymmetry_reversibility_score_results.pdf}
    \figcaption{Model reversibility scores on the Ssym dataset.}{
    Reversibility is quantified as one minus the mean squared error between the forward mutational
    effect $\Delta\log p_{fwd}$ and the negated reverse effect $-\Delta\log p_{rev}$ for different
    models (rows), where model predictions are conditioned on the protein structure containing the
    outgoing amino acid; The resulting score is normalized to lie between -1 and 1, with higher
    values indicating a greater degree of reversibility: $1 - mean((\Delta\log p_{fwd}+\Delta\log
    p_{rev})^2) / (mean({\Delta\log p_{fwd}}^2)+mean({\Delta\log p_{rev}}^2))$.
    }
    \label{fig:ssym_antisymmetry_reversibility_score_results}
\end{figure}


\begin{figure}[ht!]
    \centering
    \includegraphics[width=1.0\textwidth]{hermes/figures/heatmap.pdf}
    \figcaption{Model predictions for amino acid preferences at sites with known
    antigen-stabilizing mutations.}{ Predictions from different models (columns) are shown for sites
    with known antigen-stabilizing mutations specified in Fig.~\ref{fig:antigen_figure_main_results}
    and Table~\ref{table:antigen_results}. For each antigen (rows), predictions are computed using
    the protein structure corresponding to the PDB ID indicated on the left. Predictions are
    reported as changes in Rosetta Energy Units (REU) for Rosetta, and $\Delta \log p$ for
    ProteinMPNN and HERMES models.
    Wild-type amino acids are marked with centered crosses, while stabilizing mutant amino acids are
    indicated by dark borders. Amino acids are grouped by broad biochemical class (see legend at the
    bottom) and, within each class, ordered by increasing size (number of atoms).
    }
   \label{fig:antigen_heatmap} 
\end{figure}

\begin{figure}
    \centering
\includegraphics[width=0.95\linewidth]{hermes/figures/pvalues_vs_random.pdf}\vspace{0.5cm}
\includegraphics[width=0.95\linewidth]{hermes/figures/pvalues_vs_blosum62.pdf}
    \figcaption{Statistical significance for the number of retrieved antigen-stabilizing mutations.}{
 Shown are p-values from binomial tests comparing the number of antigen-stabilizing mutations
 retrieved by each model (rows) against random expectation (top) and the BLOSUM62 predictions
 (bottom). These significance tests correspond to the results reported in
 Figs.~\ref{fig:antigen_figure_main_results},~\ref{fig:antigen_figure_of_different_types_of_mutations}
 and Table~\ref{table:antigen_results}; see Methods for details of p-value computation.}
    \label{fig:antigen_pvalues_vs_random}
    \label{fig:antigen_pvalues_vs_blosum62}
\end{figure}


\begin{figure}
    \centering
    \includegraphics[width=0.8\linewidth]{hermes/figures/t2837_radial_plot_py_vs_bp.pdf}
    \figcaption{Comparison of PyRosetta and Biopython Pre-processing pipelines for predicting
    mutation stability effects on T2837.}{ Classification accuracy metrics, analogous to those shown
    in Fig.~\ref{fig:radial_plots__noise}, are reported for fine-tuned models (left) and zero-shot
    models (right). In each case, models trained using PyRosetta-based Pre-processing are compared
    with those using Biopython-based Pre-processing (denoted by “Bp” in the model name).
    Model labels specify the architecture, the coordinate-noise amplitude, and, when applicable, the
    fine-tuning dataset (listed after ``+"). Consistent with results on the RaSP dataset
    (Fig.~\ref{fig:rasp_exp}), Biopython-Pre-processed models show slightly reduced performance
    relative to PyRosetta-Pre-processed models; however, this difference becomes statistically
    insignificant after fine-tuning.}
    \label{fig:t2837_py_vs_bp}
\end{figure}

\begin{figure}
    \centering
    \includegraphics[width=1.0\linewidth]{hermes/figures/hermes_relaxed_ablation.pdf}
    \figcaption{Ablation of PyRosetta fastrelax parameters for HERMES-\textit{relaxed} 0.50 on
    the cDNA117k dataset.}{
    \textbf{(A)} Classification accuracy metrics, analogous to those shown in
    Fig.~\ref{fig:radial_plots__noise}, are reported for HERMES-{\em relaxed} models evaluated on
    relaxed mutant structures using different PyRosetta fastrelax parameters (indicated by color).
    HERMES-{\em relaxed} scores a mutation as the log-probability difference between the mutant and
    wild-type amino acids. The wild-type log-probability is evaluated on the wild-type structure,
    while the mutant log-probability is evaluated on the wild-type structure after introducing the
    mutation and performing local relaxation. We use the PyRosetta fastrelax protocol and vary the
    following parameters, noting that the procedure is stochastic: (1) {\bf side}, the distance
    cutoff for side-chain relaxation; (2) {\bf bb}, the distance cutoff for backbone relaxation; (3)
    {\bf nreps}, the number of protocol repetitions, with the lowest-energy conformation retained;
    (4) {\bf ens}, the ensemble size, where predictions are averaged over relaxations obtained with
    different random seeds.
    {\textbf{(B)} Inference speed for predicting mutational effects on 51 mutants across the same
    PyRosetta fastrelax parameters as in (A) (colors). A single NVIDIA A40 GPU and a single CPU with
    64G of memory were used for all parameters.}
    }
    \label{fig:relaxations_ablation}
\end{figure}


\begin{figure}
    \centering
    \includegraphics[width=1.0\textwidth]{hermes/figures/rsvf_si_figure.pdf}
    \figcaption{Structure of the wild-type pre-fusion RSV-F antigen highlighting observed and
    candidate antigen-stabilizing mutations.}{
The wild-type structure is shown on the left (PDB ID: 4JHW). A single protomer of the trimer is
displayed as a cyan cartoon, with the remaining protomers shown as a grey surface. WT denotes wild
type and MT denotes mutant. Mutations labeled as ``in silico" were introduced using PyMOL's
Mutagenesis Wizard starting from the wild-type structure. Steric clashes (orange dashed lines) were
identified using PyMOL's ``find clashes" command, and polar contacts (yellow dashed lines) were
identified using the corresponding PyMOL command.
    \textbf{(A)} The L207 mutant has been experimentally shown to stabilize the pre-fusion
    conformation~\citep{mclellan_structure-based_2013} and appears to enhance intraprotomer packing.
    We speculate that the I207 mutant, which HERMES-\textit{amortized} predicts to have a comparable
    ranking to L207 in Fig.~\ref{fig:antigen_heatmap}, would pack similarly and may therefore
    represent an additional stabilizing mutation worth screening.
\textbf{(B)} Mutant F190 is observed to be stabilizing~\citep{mclellan_structure-based_2013}, though
it appears to slightly over-pack the region. We speculate that L190, which HERMES-\textit{amortized}
predicts to have a comparable ranking to F190 in Fig.~\ref{fig:antigen_heatmap}, would provide a
similarly stabilizing effect without over-packing the region.
    }
    \label{fig:rsvf_si_figure}
\end{figure}
   
\begin{figure}
    \centering
    \includegraphics[width=1.0\textwidth]{hermes/figures/ha_si_figure.pdf}
    \figcaption{Structure of the wild-type pre-fusion Universal-HA antigen, with highlighted
    observed and candidate antigen-stabilizing mutations.}{
    The wild-type structure is shown on the left (PDB ID 7VDF).
    A single protomer of the trimer is shown as cyan cartoon; the other copies are shown as grey
    surface.
    WT denotes wild type and MT denotes mutant. Mutations labeled as ``in silico" were introduced
    using PyMOL's Mutagenesis Wizard starting from the wild-type structure. Steric clashes (orange
    dashed lines) were identified using PyMOL's ``find clashes" command, and polar contacts (yellow
    dashed lines) were identified using the corresponding PyMOL command.
    We highlight the H355W mutation which, despite stabilizing the pre-fusion conformation, exhibits
    substantial steric clashes in the wild-type structural context.
    }
    \label{fig:ha_si_figure}
\end{figure}

\begin{figure}
    \centering
    \includegraphics[width=1.0\textwidth]{hermes/figures/hmpvf_si_figure.pdf}
    \figcaption{Structure of the wild-type pre-fusion hMPV-F antigen, with highlighted observed
    and candidate antigen-stabilizing mutations.}{
    The wild-type structure is shown on the left (PDB ID 5WB0).
 A single protomer of the trimer is shown as cyan cartoon; the other copies are shown as grey
 surface.
    WT denotes wild type and MT denotes mutant. Mutations labeled as ``in silico" were introduced
    using PyMOL's Mutagenesis Wizard starting from the wild-type structure. Steric clashes (orange
    dashed lines) were identified using PyMOL's ``find clashes" command, and polar contacts (yellow
    dashed lines) were identified using the corresponding PyMOL command.
    \textbf{(A)} The proposed V449D mutation is highlighted, which likely stabilizes the complex by
    removing a surface-exposed hydrophobic residue and introducing an intraprotomer polar contact.
    We speculate that substitution with glutamic acid, which HERMES-\textit{amortized} predicts to
    have a comparable ranking to D449 in Fig.~\ref{fig:antigen_heatmap}, would produce a similar
    stabilizing effect.
     \textbf{(B)} Introduction of an Arginine in place of a Valine at position 112 likely introduces
     an intraprotomer polar contact, as well as packing against the adjacent protomer (shown in
     green); we believe L112, which HERMES-\textit{amortized} predicts to have a comparable ranking
     to R112 in Fig.~\ref{fig:antigen_heatmap}, would also pack against the adjacent protomer better
     than Valine.
    \textbf{(C)} A159L stabilizes the complex likely due to a cavity-filling effect, albeit
    necessitating nearby I137 to adopt a different rotamer; we believe I159, which
    HERMES-\textit{amortized} predicts to have a comparable ranking to L159 in
    Fig.~\ref{fig:antigen_heatmap}, would also fulfill a similar role on the condition of L141 also
    adopting a different rotamer.
    }
    \label{fig:hmpvf_si_figure}
\end{figure}
  

\begin{figure}
    \centering
    \includegraphics[width=1.0\textwidth]{hermes/figures/denve_si_figure.pdf}
    \figcaption{Structure of wild-type pre-fusion DENV-E antigen, with highlighted observed and
    candidate mutations.}{
    The wild-type structure is shown on the left (PDB ID 1OAN).
   A single protomer of the trimer is shown as cyan cartoon; the other copies are shown as grey
   surface.
    WT denotes wild type and MT denotes mutant. Mutations labeled as ``in silico" were introduced
    using PyMOL's Mutagenesis Wizard starting from the wild-type structure. Steric clashes (orange
    dashed lines) were identified using PyMOL's ``find clashes" command, and polar contacts (yellow
    dashed lines) were identified using the corresponding PyMOL command.
   \textbf{(A)} The S29K mutation is experimentally observed to be
   stabilizing~\citep{phan_conserved_2022}, despite appearing to introduce substantial steric
   clashes when inspected in the wild-type structure, suggesting that stabilization requires a shift
   in backbone conformation.
    \textbf{(B)} The F279W mutation is experimentally observed to be
    stabilizing~\citep{phan_conserved_2022}, likely by filling an under-packed cavity. A substantial
    backbone rearrangement is also observed in the mutant structure (residues 269–281), potentially
    facilitated by the T280P mutation.
    \textbf{(C)} The G106D mutation is experimentally observed to be
    stabilizing~\citep{phan_conserved_2022}, likely through the introduction of a polar contact with
    the adjacent protomer (shown in green) and/or interactions with surrounding water molecules.
    }
    \label{fig:denve_si_figure}
\end{figure}


